\section{Modulación y procesos estacionarios}\label{sec.modest}

Sea $x(t)$, proceso estacionario con autocorrelación $R_x(\tau)$,
densidad espectral de potencia $S_x(f)$. Para esta sección suponemos también que
el valor medio $E[x(t)]=0$.

Nos interesa estudiar la
señal modulada $x(t)\cos(\omega_ct)$ con los mismas técnicas. Sin
embargo, este proceso no es estacionario, ya que se anula en
momentos fijos del tiempo. ?`Cómo proceder entonces?
Estudiamos dos alternativas.

%\subsection*{1. Introducir una fase aleatoria}
\begin{enumerate}
\item Modular con fase aleatoria.

Como ya se vio, una sinusoide de fase aleatoria uniforme es un proceso estacionario. Definimos entonces
\[z(t)=x(t)\cos(\omega_ct+\theta),\qquad \theta\sim U[0,2\pi],\text{ indep de }x(t).\]
Usando la independencia, se ve que $z(t)$ es un proceso estacionario cuya autocorrelación es el
producto de las autocorrelaciones de $x(t)$ y del coseno de fase
aleatoria:
\begin{align*}
    R_z(\tau)&=E\left[x(t)\cos(\omega_ct+\theta)x(t-\tau)\cos(\omega_ct+\theta-\omega_c\tau)\right]
    \\ & = R_x(\tau)E\left[\cos(\omega_ct+\theta)\cos(\omega_ct+\theta-\omega_c\tau)\right]
    \\ &=\frac{1}{2}R_x(\tau)\cos(\omega_c\tau).
\end{align*}
La correspondiente densidad espectral de potencia es
\[
 S_z(f) =\frac{S_x(f-f_c)+S_x(f+f_c)}{4}.
\]
Este modelo puede ser útil para estudiar situaciones en que no hay sincronismo entre la portadora y el proceso que modula.

\item Introducir una componente en cuadratura.

Complementamos aquí $x(t)$ con otro proceso $y(t)$ ``en
cuadratura",  y consideramos el proceso con envolvente compleja
$x(t) + j y(t)$, es decir
\[
u(t)=x(t)\cos(\omega_ct)  -y(t)\sen(\omega_ct).\]

Suponemos que $y(t)$ es un proceso aleatorio estacionario de iguales características que
$x(t)$, en particular de media 0 y $R_y(\tau)=R_x(\tau)$. Suponemos inicialmente que ambos procesos son
\emph{incoherentes}, es decir
\[
R_{xy} (\tau) = E[x(t)y(t-\tau)] \equiv 0.\]
En ese caso tendremos
\begin{align*}
                E\left[u(t)u(t-\tau)\right]  = &
                \ R_x(\tau)\cos(\omega_ct)\cos\left(\omega_c(t-\tau)\right)
                                 + R_y(\tau)\sen(\omega_ct)\sen\left(\omega_c(t-\tau)\right) \\
                                 = & \ R_x(\tau)\cos(\omega_c\tau),
            \end{align*}
en el que se utilizó para el último paso una identidad trigonométrica. Concluimos entonces que
el proceso $u(t)$ es estacionario, con densidad espectral de potencia
\[
S_u(f)=\frac{S_x(f-f_c)+S_x(f+f_c)}{2}.
\]
El hecho que $S_u(f)=2 S_z(f)$ se explica intuitivamente porque a través del proceso $y(t)$ se introdujo
``el doble de ruido".

De hecho, el procedimiento anterior puede generalizarse al caso en que $y(t)$ tenga correlación cruzada no nula
$R_{xy}(\tau)$, bajo cierta restricción, lo cual se aplica al estudio del ruido pasabanda que se verá a continuación.
Como esa generalización es algo técnica la incluimos en el Apéndice A al final de estas notas.
\end{enumerate}
